\section{系统组成及功能说明}
\subsection{整体介绍}
系统由图像数据平面和软件控制平面组成。数据平面连接摄像头、RAW处理、DDR、CNN和HDMI；控制平面由Sapphire RISC-V、APB寄存器、按键和手势状态机构成。\sig{video_system_720}在100MHz UI域组织帧池、作业与DDR仲裁，通过异步FIFO与166.667MHz CNN域、74.25MHz显示域交换数据。摄像头字节时钟来自D-PHY，可能在低功耗期间停止，因此不直接跨域采样多位包头。

\begin{table}[H]\centering\small
\caption{模块与主要源码对应}
\begin{tabularx}{\linewidth}{@{}p{.27\linewidth}X@{}}\toprule 功能 & 主要源码\\\midrule
PPI与CSI解析 & \sig{csi_ppi_bridge_recover.sv}\\
RAW预处理 & \sig{camera_bayer_720.sv}\newline \sig{camera_pack_720.sv}\\
源帧与作业管理 & \sig{camera_frame_writer_720.sv}\newline \sig{source_frame_pool.sv}\newline \sig{video_job_control.sv}\\
分条DMA与跨域 & \sig{video_tile_engine_async.sv}\newline \sig{video_cnn_async.sv}\\
卷积及参数 & \sig{full_dw3x3.sv}\newline \sig{full_pointwise.sv}\newline \sig{full_parameter_bank.sv}\\
显示及UI & \sig{frame_reader.sv}\newline \sig{video_raster60.sv}\newline \sig{fps_overlay60.sv}\newline \sig{tmds_hdmi60.sv}\\
CPU及交互 & \sig{video_control_apb_fast.sv}\newline \sig{style_keys47.sv}\newline \sig{gesture_snapshot.sv}和C固件\\\bottomrule
\end{tabularx}
\end{table}
除特别注明外，摄像头解析、帧池、作业、DDR用户接口及APB均使用100MHz的clk和高有效reset；CNN使用166.667MHz的\sig{fast_clk}，扫描使用74.25MHz像素时钟。本文将同组总线适当合并列出，低电平命名的外部复位在顶层转换，不与内部reset混用。

以下接口表列出主要模块边界，不重复展开全部诊断端口及厂商DDR PHY。方向均相对该节模块；有效/就绪型流仅在valid与ready同时有效时前进，无ready的摄像头输入必须持续接收。

\subsection{硬件系统介绍}
\subsubsection{平台与电气接口}
系统采用Ti60F225I3 DemoBoard V4及demo10配套SC431HAI排线模组，沿用板载DDR3、USB串口和HDMI，未增加自制机械或供电结构。板级依据为20261012版原理图\cite{board}；采集参考V7，对照V9的Bayer与数据延迟处理\cite{demo}，未改原demo。
\begin{table}[H]\centering\small
\caption{外设连接与电平约束}
\begin{tabularx}{\linewidth}{@{}p{.22\linewidth}p{.18\linewidth}X@{}}\toprule 接口 & 配置电平 & 连接与实现约束\\\midrule
MIPI专用RX & 1.2V & Bank2A/2B；配套模组接J4，J2/J3均2--3\\
摄像头控制 & 3.3V & SCL D12、SDA B14开漏；C12模组RESET高有效\\
四按键 & 3.3V输入 & KEY0--3依次C4、P14、N14、A3，低有效\\
板载USB串口 & 3.3V & R4 RX、R14 TX；J8不是UART接线座\\
DDR3 & 1.5V & 保留原板DDR PHY、终端及OE绑定\\
HDMI串化侧 & 1.8V & Bank3A，使用现有板级TMDS电路\\\bottomrule
\end{tabularx}
\end{table}
电气签核逐项检查79个实际端口、12组Bank、方向、I/O标准及DDR数据/OE绑定，与CAD时序检查相互独立。用户已确认J4和跳线接法；实际电源电压、连接器状态和温升仍须实板测量，数字审查不代替物理测量。


\blocktitle{板级连接与内部总线的分工}
图\ref{fig:board-direction}区分电气接口与片内数据接口。DDR引脚侧为16位双沿数据总线，用户侧则是100MHz、128位AXI，二者通过厂商控制器转换；不能把AXI位宽或UI频率直接套到外部DQ上。MIPI只承担输入，HDMI只承担输出，GPIO控制信号与高速差分信号分别使用对应Bank供电。
\begin{figure}[H]\centering
\begin{tikzpicture}[x=1cm,y=1cm,every node/.style={font=\small,align=center},
box/.style={draw=navy,rounded corners=2pt,minimum height=.85cm,text width=3.1cm},
arr/.style={-{Stealth},thick,draw=teal}]
\node[box] (cam) at (0,1.3) {J4摄像头模组\\MIPI RX / 控制GPIO};
\node[box] (key) at (0,-1.3) {四按键 / 板载USB\\输入 / 串口诊断};
\node[box,minimum height=3.6cm,text width=3.3cm] (fpga) at (5,0) {Ti60F225I3\\[4pt]D-PHY / DDR PHY\\RISC-V / CNN\\TMDS编码及串化};
\node[box] (ddr) at (10,1.3) {U12 DDR3\\DQ / DQS双向};
\node[box] (hdmi) at (10,-1.3) {HDMI连接器\\TMDS输出};
\draw[arr] (cam.east)--node[above,font=\scriptsize]{RAW10}(fpga.west|-cam.east);
\draw[arr] (fpga.west|-0,.6)--node[below,font=\scriptsize]{I²C / RESET}(0,.6)--(cam.south);
\draw[arr] ([yshift=2mm]key.east)--node[above,font=\scriptsize]{按键输入}([yshift=2mm]fpga.west|-key.east);
\draw[arr] ([yshift=-2mm]fpga.west|-key.east)--node[below,font=\scriptsize]{UART诊断}([yshift=-2mm]key.east);
\draw[arr,<->] (fpga.east|-ddr.west)--node[above,font=\scriptsize]{数据 / 选通}(ddr.west);
\draw[arr] (fpga.east|-hdmi.west)--node[above,font=\scriptsize]{视频 / 时钟}(hdmi.west);
\end{tikzpicture}
\caption{板级接口方向示意，省略电源和所有单根管脚。配置电平以表格和79端口审查清单为准。}
\label{fig:board-direction}
\end{figure}
摄像头I\textsuperscript{2}C不是推挽高低输出。固件写入\sig{scl_release/sda_release}，置1表示高阻释放，由模组接口上拉形成高电平；置0才驱动低电平，并从输入端读回总线状态。板上配套模组的3.3V控制兼容性不能外推到裸传感器。复位信号的内部名称含\sig{sensor_reset_n}，但连接模组时须结合顶层极性变换判断，不能单凭后缀认定引脚低有效。

DDR必须在PLL锁定、PHY复位释放和校准完成后接收图像访问；校准失效时撤销源帧有效性。HDMI发送使用板上既有耦合、偏置和保护电路，保持原配置与方向。

\subsubsection{摄像头接收与RAW预处理}
\blocktitle{作用}
将四lane CSI-2 RAW10包恢复为可发布的RGB565帧。\sig{csi_ppi_bridge_recover}将字节及SoT等状态形成44位记录，经512深异步FIFO送入UI域，保持数据与状态原子对应。解析器检查包头、ECC/CRC、字节数与帧长度，错误帧不发布。
\begin{porttable}
四lane PPI字节及状态 & 输入 & 32位数据与各lane有效、同步及SoT，来自D-PHY RX\\
\sig{vs/hs/valid/raw} & 输出 & 帧、行、样本有效及40位RAW10四样本\\
\sig{packet_error/overflow} & 输出 & 包错误与FIFO溢出，影响整帧有效性\\
\sig{frame_start/end/bad} & 输出 & 预处理后边界与坏帧标志，供打包及手势旁路\\
\sig{pixels[31:0]} & 输出 & 两个RGB565像素，低16位为较早的左像素\\
\end{porttable}
传感器配置为1920×1080四lane RAW10，每行2400字节，每五字节恢复四个10位样本。Bayer模块以邻行、邻列插值补齐颜色，再按约$x=\lfloor3u/2\rfloor$、$y=\lfloor3v/2\rfloor$映射形成1280×720全视野输出，RAW高八位量化为RGB565。尺寸变换在CNN之前进行；风格网络内部不再减半分辨率。帧尾排空插值与打包流水，核对像素数，避免漏掉最后一行。


\blocktitle{协议解析状态与错误边界}
接收器分为字节域记录、UI域lane队列和协议状态机三层。每条lane的小队列消化到达差异，只有需要的lane数据齐备才弹出对应字节。其主要状态为\sig{WAIT_SOT}、\sig{HEADER}、\sig{PAYLOAD}、\sig{CRC}：等待同步、识别4字节包头、累计指定长度负载，再比较CRC。短包携带帧边界，长包携带RAW行。当前范围是VC0、四lane、每次HS传输一个CSI包，RAW10长包datatype为\sig{0x2B}，行word count固定2400。

ECC检查采用检测后拒绝策略，包括可纠正单比特错误也不自动修复。CRC错误、缺行、长度不符或包中断均使当前帧失效。仅在所有lane队列已空、处于等待SoT且该记录不含数据或同步时，将孤立SoT状态记为软事件；这不会跳过随后的帧边界、ECC、CRC和尺寸检查。异步FIFO溢出属于保持故障，需要复位恢复。由此将“传输恢复”与“接受一幅好图”分开。

\blocktitle{RAW10解包、颜色与位置一致性}
每组负载字节记为$b_0,\ldots,b_4$，四个RAW样本恢复为
\[
p_i=(b_i\ll2)\;|\;((b_4\gg2i)\,\&\,3),\quad i=0,1,2,3.
\]
输出\sig{raw[39:0]}由低到高放置$p_0$到$p_3$。预处理取各样本高八位，两拍处理四个RAW样本；两条行缓存与横向历史组成3×3窗口。以BGGR为例，蓝色位置保留中心B，四个正交邻点平均得到G，四个对角邻点平均得到R；绿色位置依据行列奇偶选择水平或垂直邻点恢复R/B。Bayer模式在帧起点锁存，邻域数据、行列奇偶与有效标志一起延迟。

尺寸变换保留每组三个RAW坐标中的前两个，对应源坐标序列$0,1,3,4,\ldots$。最后输出坐标映射到RAW的1918列、1078行，仍可取得右、下邻点；左、上边界使用reflect101。它是插值后的取样，不是通用抗混叠缩放滤波器，也不是把640×360重复放大。图像偏绿需要分别检查Bayer相位、通道排列及照明白平衡，单改风格权重不能修复输入颜色错误。

\begin{table}[H]\centering\small
\caption{相邻模块的数据布局契约}
\begin{tabularx}{\linewidth}{@{}p{.24\linewidth}p{.21\linewidth}X@{}}\toprule
边界 & 位宽/布局 & 推进与帧尾条件\\\midrule
D-PHY→PPI FIFO & 44位：状态12＋数据32 & 数据与状态同一记录跨域\\
CSI→Bayer & 40位，4个RAW10 & \sig{valid}限定样本，\sig{vs/hs}限定帧行\\
Bayer→打包 & 32位，2个RGB565 & \sig{pixel_valid}一次推进2像素\\
打包→DDR writer & 179位记录 & 独立EOF；不以最后像素代替好帧结束\\
DMA→CNN & 24位，R在低8位 & valid/ready同拍握手才消费\\
显示FIFO→raster & 16位RGB565 & 有效区逐像素读取，无数据则计欠载\\\bottomrule
\end{tabularx}
\end{table}

\subsubsection{DDR写入与源帧池}
\blocktitle{作用}
将无背压的摄像头像素流与受仲裁影响的DDR分离。128位字由低至高排列8个RGB565像素，行stride为2560B，每帧1843200B，即115200个128位字。写入记录使用17位字索引，避免超过65535字后回绕。
\begin{porttable}
\sig{fifo_data[178:0]} & 输入 & 含帧号、字索引、坏帧及独立EOF的像素记录\\
\sig{allocate/allocation_ack} & 双向 & 写帧前申请空闲bank并获得编号\\
\sig{AW/W/B} & 双向 & 写地址、128位数据、最终响应；B错误作废本帧\\
\sig{publish/abort_capture} & 输出 & 发布完整好帧或丢弃，半帧不能成为有效源\\
\sig{claim/release} & 双向 & 引擎领取并锁住源bank，使用后释放\\
\end{porttable}
\sig{source_frame_pool}管理4、6、8MiB三个源bank，轮流承担采集、最新完整帧及引擎锁定帧。一个bank不能同时被覆盖和读取。推理落后时允许新好帧替换未领取旧帧，并显式累计跳过数；持续采集不等于每帧都已经推理。


\blocktitle{DDR中的地址安排}
所有bank以2MiB为槽跨度，一幅720p RGB565占用约1.758MiB，剩余空间不作为下一帧有效像素。行主序、无行间padding，坐标$(x,y)$对应字节地址
\[
A(x,y)=A_{\rm bank}+2560y+2x.
\]
\begin{table}[H]\centering\small
\caption{六个图像槽的地址及所有权}
\begin{tabularx}{\linewidth}{@{}p{.23\linewidth}p{.31\linewidth}X@{}}\toprule
图像槽 & 基地址 & 谁能写、何时释放\\\midrule
输出0 / 1 / 2 & \sig{0x000000 / 0x200000 / 0xA00000} & 引擎写；前台槽需显示commit ACK后释放\\
原图0 / 1 / 2 & \sig{0x400000 / 0x600000 / 0x800000} & 摄像头写；CNN领取后锁至作业结束\\\bottomrule
\end{tabularx}
\end{table}
打包记录的位定义为\sig{[178]EOF}、\sig{[177]bad}、\sig{[176:145]frame_id}、\sig{[144:128]word_index}和\sig{[127:0]pixels}。四拍双像素合成一个128位词，先来的像素处于低位。EOF是独立记录，其index应等于115200；它证明帧边界已到达，而不是只证明缓存中曾出现足够像素。

\blocktitle{写入状态机和发布屏障}
\sig{camera_frame_writer_720}经过IDLE、ALLOCATE、ALLOC\_WAIT、WRITE和RESP。检测到新帧首词先申请bank，收到allocation ACK再开始写；WRITE分别记录AW和W是否握手，允许二者不同拍完成。进入RESP后等待真正B响应，成功后才读取下一记录。EOF只有在精确词数、连续index、同frame\_id、无bad且此前B全部成功时触发publish。写数据握手完成不能提前当作帧可用。

源池用\sig{capture_owned}、\sig{engine_owned}与\sig{latest_valid}表述所有权。空闲槽的选择同时排除三个占用角色。若claim与新publish同拍，旧latest交给引擎，新好帧成为latest；引擎仍读取旧帧，采集不会覆写它。未领取的latest被下一好帧替代时累加skipped计数。分配失败或FIFO积压引起的坏帧与这种主动舍弃必须分开统计。

这种安排允许推理速度暂时落后，同时保持读取帧完整。代价是系统倾向处理最新可用帧，不能承诺输入每一帧都被推理；实际是否发生跳帧需读取运行计数。

\subsubsection{作业控制与输出槽}
\sig{video_job_control}负责源帧及输出槽所有权，APB与系统集成逻辑锁存作业风格和视图。忙中仅保留最新用户请求，不改变正在计算的身份。三个输出槽位于0、2、10MiB，轮流承担计算、待提交和前台，最多维护两项完整输出队列。
\begin{porttable}
\sig{start/raw_mode} & 输入 & 作业启动与原画标志，上层先检查源帧与权重\\
\sig{claim/claim_ack} & 双向 & 领取完整源帧，并返回锁定的source\_frame\\
\sig{engine_start/done/error} & 双向 & 启动引擎；done晚于最后DDR B成功响应\\
\sig{output_bank} & 输出 & 锁定的目的槽；源地址由源池bank换算\\
\sig{commit/commit_ack} & 双向 & 安全消隐提交与显示确认，ACK后释放旧前台\\
\end{porttable}
左原右风格在CNN内已经对齐的原RGB支路上，以全局$x<640$选择原图，右侧选择风格结果，二者来自同一源帧与坐标；不将“最新原图”叠到旧风格帧。每个输出槽保存视图及风格元数据，前台身份只在commit ACK后更新。


\blocktitle{计算完成与显示完成为什么分开}
计算完成表示全部输出已写入DDR，显示完成表示显示端已经安全采用这一bank。若把两者串行绑定，25.6ms的计算还可能等待下一次16.67ms显示边界。当前控制器在计算结束后立即释放源帧并将结果排队，下一作业可在另一个空闲输出槽启动，显示端独立消费完整队列。

作业状态机为IDLE→CLAIM→RUN。IDLE检查空闲输出槽和队列容量；CLAIM等待源池回传锁定的bank/frame\_id；RUN发出一次engine\_start并等待done或error。error优先于同拍done，失败作业释放源但不入显示队列。队列最多两项，且旧前台占用独立保存；显示reader提前改变观察用的front\_bank并不能释放旧槽。

\begin{figure}[H]\centering
\begin{tikzpicture}[every node/.style={font=\small,align=center},n/.style={draw=navy,rounded corners,text width=2.7cm,minimum height=1cm},a/.style={-{Stealth},thick,teal}]
\node[n] (a) at (0,0){空闲输出槽\\领取完整源帧};
\node[n] (b) at (3.9,0){RUN\\计算并写DDR};
\node[n] (c) at (7.8,0){READY队列\\最后B已成功};
\node[n] (d) at (11.7,0){FRONT\\显示安全ACK};
\draw[a] (a)--(b);\draw[a] (b)--(c);\draw[a] (c)--(d);
\draw[a] (d.south)--++(0,-.7)-|node[pos=.25,below]{下一次换帧ACK后，旧前台才回收}(a.south);
\end{tikzpicture}
\caption{输出槽生命周期。源帧锁在计算结束时释放，旧显示槽锁在下一次安全换帧ACK时释放。}
\label{fig:ownership}
\end{figure}
对比模式在计算结果形成时合成，不需要显示端额外再读一幅原图。原RGB由独立FIFO保存到对应CNN增量到达，两支路都有效才接受下游ready；因此背压不会把第$n$个原像素与第$n+1$个风格像素相加。分条坐标加上origin恢复全局位置，左半边原图与右半边风格的分界始终是全图第640列。

\subsubsection{分条DMA与总线仲裁}
\sig{video_tile_engine_async}将720p拆为四条：每条有效320列，左右各16列halo，实际输入352×720。四条各输出320×720，拼回完整921600像素。halo只提供邻域，不作为重复像素发布；边界按全局图像坐标作reflect101处理。
\begin{porttable}
\sig{clk/fast_clk/reset} & 输入 & UI100MHz与CNN166.667MHz，原子FIFO及握手跨域\\
\sig{start/style/raw_mode} & 输入 & 作业启动与模式标志，启动时锁存\\
\sig{cfg_valid/address/data} & 输入 & 配置写入；由cfg\_ready反馈接收状态\\
\sig{AR/R}\newline \sig{AW/W/B} & 双向 & 30位字节地址、128位数据，核对响应和RLAST\\
\sig{busy/done/error/cycles} & 输出 & 作业状态与UI域周期计数\\
\end{porttable}
引擎采用突发读，两组128位写打包槽重叠拼包和AXI事务。读仲裁优先响应同时请求的显示读取，授予事务保持至RLAST。写信用由真实B响应释放，不以W通道完成代替DDR写完成。背压期间保持地址、数据与valid，异常响应不先写入有效显示FIFO再报错。


\blocktitle{四条带的几何关系}
第$k$条有效列为$320k\ldots320k+319$，输入原点为$320k-16$。四条输入横坐标范围分别是$[-16,335]$、$[304,655]$、$[624,975]$和$[944,1295]$。只保留本条局部$x=16\ldots335$的结果。31×31感受野的半径是15，16列halo覆盖邻域并留出对齐余量。

每一层按全局图像边界处理reflect101，不能将条带接缝误当作独立图像边缘。条外填充位置不作为最终输出；极左/右边界的镜像取样由算子层的全局有效区域逻辑完成。逻辑计算像素量为$4\times352\times720=1013760$，比完整图多10\%；实际DDR读流量还受256B缓存块对齐及重取影响，不能只按10\%估算总线开销。

\blocktitle{读缓存、写打包与错误排空}
读地址按256B对齐，每个突发16拍、每拍128位，装入16词局部缓存；按像素$x$的高位选缓存块、低位选词和16位lane。单个突发不跨4KiB边界。显示与CNN使用读仲裁器，授予后维持所有权至RLAST，避免两主机混收一条响应。

输出按8个RGB565像素合成一词，两个打包槽轮流承担“正在发送”和“正在拼下一词”。只有B响应退休该槽；W完成但B未回时槽仍占用。引擎一笔写与摄像头一笔写可由\sig{axi_write_pair}形成总共两笔在途，所有者队列按固定ID响应顺序路由B。它不是无限发射的多ID写DMA。

遇到RRESP错误、RLAST位置错误或BRESP错误，停止接受新的有效输出，先排空已获总线接受的事务，随后报告error。四条全部结束、保留像素数等于921600、写队列为空且最终B成功才发done；任何单条done都不能发布半幅图。

\blocktitle{带宽和算力分别预算}
仅按有效RGB565字节，30fps采集、30fps推理读写及60Hz显示至少需要
\[
1.8432\,{\rm MB}\times(30+30+30+60)=276.48\,{\rm MB/s}.
\]
该式不含halo、缓存重复读取、DDR刷新和仲裁空洞。128位100MHz UI的每方向理论数据峰值为1.6GB/s，但单拍写响应、行切换和共享存储器会显著降低有效率，必须用真实并发计数验证。卷积名义运算量为396 MAC/像素，包含halo约4.0145亿MAC/帧，30fps约120.44亿MAC/s；这解释了为何需要流水、并行和片内复用，不能只看权重文件很小便认定容易实时运行。

\subsubsection{CNN卷积与参数存储}
\sig{full_dw3x3}执行六通道独立3×3膨胀卷积，多行缓存保存历史，列地址分为三个bank，内部位置并行读取三tap。反射边界遇到bank冲突时用第二读阶段补齐，再进入乘加，不复制整份历史以增加端口。\sig{full_pointwise}执行1×1投影、通道混合与有符号输出头，采用两输出通道lane、四像素上下文重叠计算。context标签随流水传递，按接收顺序退休，保持像素顺序。
\begin{porttable}
\sig{in_data/valid/ready} & 双向 & DW输入48位六通道uint8；PW依CIN参数确定\\
\sig{in_skip/in_inside} & 输入 & 同像素残差及全局有效区域标志\\
\sig{out_data/valid/ready} & 双向 & 中间六通道uint8；head为三个有符号16位值\\
\sig{origin_x/start} & 输入 & 条带全局横向原点和初始化，决定图像边界\\
\sig{cfg_valid/address/data} & 输入 & 权重、偏置、逐通道移位的原子配置\\
\end{porttable}
\sig{full_parameter_bank}按layer、kind、item选中参数；APB控制器检查范围与装载顺序。十组风格只保留一个活动工作集，共用计算核。增加风格主要增加ROM和目录，仍有存储与控制开销，不能说风格数完全不影响资源。


\blocktitle{DW：用行历史换取空间邻域}
DW对每个通道独立计算$z_c(x,y)=\sum_{i,j}w_{c,i,j}h_c(x+di,y+dj)+b_c$，$i,j\in\{-1,0,1\}$。历史RAM按条带宽度配置，只保存计算窗口所需的行，避免缓存整幅六通道特征图。三个列bank允许一个阶段读取一组三个横向tap，乘加结果再按行组归并；边界镜像发生端口冲突时补读而不产生错误重复值。

窗口形成、地址计算、同步RAM读出、乘积和加法树分别寄存，局部行基址提前计算，缩短“坐标→bank→RAM地址→乘加”的组合链。每次新条带start重置有效历史及坐标，不能借用上一条残留。RAM旧值存在不等于有效窗口存在，必须同时检查inside和行列有效条件。

\blocktitle{PW：通道维复用与有序输出}
PW计算$z_o=\sum_{c=0}^{C_{\rm in}-1}w_{o,c}h_c+b_o$。两条输出lane并行处理两个输出通道，四组context保存不同像素的输入、中间累加和完成状态，在乘加等待期间处理其他context。结果按输入顺序退出；不能按哪个context先算完就先输出，否则后续行缓存的空间坐标会错位。输入端ready只在存在可用context且下游容量允许时成立。

权重与结果写端使用固定生成的byte/channel写者，避免二维数组的动态局部写入在综合映射中产生差异。曾出现RTL通过但实际原语结果错误的历史问题，因此当前验证包含映射后EFX算子的数值核对，而非仅做RTL仿真。

\blocktitle{残差、激活与参数责任边界}
每个DW/PW残差块在PW输出与原块输入对齐后相加，饱和到uint8；输出头产生有符号RGB增量，允许减暗、加亮及改变颜色。只允许正增量会损害水墨等风格，因此头部不能套用中间层的无符号激活。BN在离线折叠，片上执行的是整数乘加、移位、饱和及残差。

\sig{full_parameter_bank}负责将合法配置写到对应寄存器，其本身不负责完整装载校验。合法顺序、数量、忙闲互斥及提交身份由APB控制器检查；跨域桥保证一次配置数据稳定到CNN确认。参数包形成、总线检查和存储写入分别有明确责任，避免以参数bank存在就误认为加载事务已经安全。

\subsubsection{HDMI与统计叠加}
\sig{frame_reader}读取已提交输出槽，经FIFO送入像素域。\sig{video_raster60}以74.25MHz产生1650×750总周期、1280×720有效区的720p60时序；RGB、DE和同步按一致延迟送入TMDS编码与串化。
\begin{porttable}
\sig{front_base/commit} & 输入 & 完整前台地址与换帧请求，只在安全窗口采纳\\
\sig{AR/R} & 双向 & DDR显示读取，检查响应后才接收像素\\
\sig{rgb/hsync/vsync/de} & 输出 & RGB888、同步及有效区，像素域对齐\\
\sig{captured/inferred} & 输入 & 好采集帧及神经作业完成计数，按UI时基取差\\
\sig{tmds_data/clock} & 输出 & 由现有接口资源串化至板级TMDS通道\\
\end{porttable}
\sig{fps_snapshot}以稳定toggle邮箱将C/I、手势进度和风格送到像素域，收到ACK前不改变数据。I仅统计完成的神经作业，原图拷贝和重复刷新不计推理FPS。KEY3关闭文字和青框，不改变RAW/SPLIT的识别资格。


\begin{table}[H]\centering\small
\caption{720p60扫描参数，正极性HS与VS}
\begin{tabular}{@{}lrrrrr@{}}\toprule
方向 & 有效 & 前肩 & 同步 & 后肩 & 总计\\\midrule
水平（像素时钟） &1280&110&40&220&1650\\
垂直（行） &720&5&5&20&750\\\bottomrule
\end{tabular}
\end{table}
刷新率由$74.25{\rm MHz}/(1650\times750)=60{\rm Hz}$确定。计数器以同步区开始，故有效区起点为$(H_{\rm sync}+H_{\rm back},V_{\rm sync}+V_{\rm back})=(260,25)$；这与从有效区开始计数只是原点不同，不能将两种坐标直接混用。RGB565扩展到RGB888时复制高位补低位，HS/VS/DE和RGB在同一像素级保持对应关系。

显示reader的换帧过程为：检测消隐toggle请求；结束旧AR/R突发；在安全窗口选择新bank；复位像素FIFO至少16个UI拍；等待两域reset busy退出；再次确认仍在安全窗口后发送ACK。窗口在有效区前一整行关闭，迟到的请求保留到下一消隐。旧front只有收到这个晚ACK才能被写方复用。

第一幅完整图尚未提交时显示彩条；有效前台存在但有效区数据缺失时，本帧标为欠载并显示诊断色，计数每帧只加一次。这样可以区分“还没有好源帧”和“显示读取供不上”。屏幕有60Hz时序并不证明CNN每秒产生60幅新图。UI中的C/I采用计数差，诊断串口用于进一步区分摄像头错误、推理完成和显示欠载。

\subsubsection{时钟复位与CDC}
各域在时钟稳定后同步释放复位；PLL失锁或DDR校准失效使数据流退出运行，不继续信任旧源帧。单比特用同步器/toggle，像素流用异步FIFO，多位配置和状态用稳定束握手。
最终约束保留23组标量首级例外、166条Gray最大延迟约束、87位稳定数据束及8条精确D-PHY首级例外。八个D-PHY首级仅连接对应二级，最小级间setup余量+13.391ns。未用整组异步时钟切断隐藏关系；旧CSI5.15元数据争议路径已不在当前公开PPI解析方案中。


\begin{table}[H]\centering\small
\caption{主要跨域通道及保持规则}
\begin{tabularx}{\linewidth}{@{}p{.25\linewidth}p{.24\linewidth}X@{}}\toprule
通道 & 机制 & 不变量\\\midrule
byte→UI PPI & 异步FIFO & 包数据与同步/错误位为一个原子记录\\
UI↔CNN像素 & 异步FIFO & ready为流量反馈，不能跨域直传单拍valid\\
UI→CNN配置/启动 & 稳定数据束＋握手 & ACK前地址、值、风格及起点保持\\
pixel↔UI换帧 & toggle＋同步器 & 请求未确认不重复翻转，受安全窗口约束\\
UI→pixel统计 & 快照邮箱 & C/I与UI状态同批锁存，ACK前不改写\\
按键/I²C回读 & 两级同步 & 异步信号只进入首级，业务逻辑取二级\\\bottomrule
\end{tabularx}
\end{table}
对异步FIFO，Gray指针的多位变化时间也必须约束；仅声明两个时钟无关会丢失这项检查。对稳定束，数据先锁住再发送toggle，接收方同步请求后取数并回ACK，约束仍限制数据传播。DPHY首级例外只落在已证明唯一二级后继的寄存器上。字节时钟在LP期间停振时，UI仍能处理超时和复位；不能依赖停止的时钟发回一次完成脉冲。

最终setup最小余量0.108ns表明目标时钟下CAD通过，但余量有限。该结果绑定当前器件速度等级、工具、布局种子、源码和约束；改动CNN位宽、参数、管线或频率后均需重做对应签核。硬件故障风险、数字CDC正确性和算法画质属于不同问题，分别通过电气审查、时序/仿真及实物图像检查。

\subsection{软件系统介绍}
\subsubsection{RISC-V固件与APB控制}
Sapphire RV32IM程序运行于32KB片内RAM，当前固件10588B，BSS结束于\sig{0xF90071AC}，栈余3668B。CPU负责控制，整幅像素由DMA经AXI访问DDR。
固件先检查处理器能力，经开漏I\textsuperscript{2}C读取\sig{0x3107/0x3108}，确认预期ID \sig{CD6B}，再写摄像头配置并读回尺寸、VTS等。无应答、ID错误或配置失败时关闭采集、复位模组、释放I\textsuperscript{2}C并禁止CNN启动。
\begin{table}[H]\centering\small
\caption{APB主控制接口，基址\sig{0xF8100000}}
\begin{tabularx}{\linewidth}{@{}p{.25\linewidth}X@{}}\toprule 偏移 & 用途\\\midrule
\sig{0x04/0x08} & 采集/复位/UI控制及busy、源帧、权重状态\\
\sig{0x0C/0x10} & sticky中断W1C及中断屏蔽\\
\sig{0x28/0x2C/0x30} & 请求风格、加载身份、clear/commit/START\\
\sig{0x50/0x54/0x58} & 配置地址、数据和确认事务数\\
\sig{0x80--0xBC} & 手势快照状态、seq/epoch及原子风格请求\\
\sig{0x2000--0x3FFC} & 2048个32位RGB565双样本词，只读\\
\sig{0x4000--0x5CBF} & 十组学习参数ROM，10×736B\\\bottomrule
\end{tabularx}
\end{table}
每组由CPU发出216项配置，等待CNN接收后计数，地址、事务与身份都正确才commit。加载目标与后来按键请求分开保存，不能把尚未加载的请求当成当前权重。引擎结束通过PLIC IRQ16通知CPU，启动及完成与完整帧所有权关联。


\blocktitle{启动顺序与失败出口}
复位后先运行片内Boot RAM中的程序，设置栈、清零BSS并配置机器中断入口，检查M扩展及乘除运算。随后释放I²C线、按模组要求控制复位，读回器件ID，再执行待机寄存器表；尺寸、VTS等读回符合预期后才开流。传感器模拟与PLL寄存器沿用已核对例程，720p由FPGA预处理形成，不凭经验修改为未经验证的传感器时序。

任何ID、ACK、总线超时或配置检查失败都进入\sig{fail(reason)}：保存错误码，关闭capture、保持模组复位、释放SCL/SDA，并停留等待人工处理。故障路径不反复开流，不以“已经有PLL时钟”当作初始化成功。I²C的START、地址高/低字节、数据和STOP由C逐步操作，SCL释放后读回实际高电平，超时就退出，避免外设拉住总线时无限等待。

\blocktitle{模型装载的完整事务}
CPU先把目标槽写入\sig{0x2C}并向\sig{0x30}写2开始清除旧装载状态；从\sig{0x4000+736s}读对应包，依次写十层的权重、偏置、移位。地址字段为layer[15:12]、kind[11:10]和item[9:0]，kind 0/1/2分别表示权重、偏置和移位。权重一次写四个INT8，item按字节索引递增4；偏置和移位每通道各写一次。

APB写\sig{0x54}时如跨域配置仍未确认，PREADY拉低延长ACCESS阶段。CPU不能把一次write指令返回前的中间状态当成CNN已写入；配置桥在实际CNN寄存器写入边沿后才ACK。当前APB校验顺序、移位不超过20、事务计数及busy互斥，当前偏置由导出器进一步证明满足22位安全范围。在线接受任意外部包时还需重新验证这个数值上界，不能仅凭现有地址检查放行。

十层全部完成后，CPU确认\sig{0x58}=216，再向\sig{0x30}写4提交；weights\_valid置位且loaded\_style对应真正装载目标。最新requested\_style与load\_target分开，加载期间用户再按键只形成下一请求。只有空闲、源帧有效、权重有效才向\sig{0x30}写1启动，避免半组新旧权重混用。

\blocktitle{中断服务与主循环}
完成和失败状态保存在sticky中断寄存器，PLIC IRQ16通知CPU。中断函数先claim，读取并缓存pending，再按读取位写1清除，最后complete同一IRQ；中断服务只处理状态，不执行整帧像素拷贝。主循环先调度符合条件的视频作业，再服务手势快照和看门狗，因此即使引擎空闲或源帧暂时停止，手势状态也不会因为仅在busy分支运行而卡死。

CPU内存同时容纳程序、全局变量、背景、特征、模板和栈。四个模板×十类×192字节为7680B，占用明显；32KB RAM剩余栈约3668B，故前景连通分析复用临时数组而非另申请一幅大图。不能在此版本直接加入大型神经分类器而忽略链接布局。

\subsubsection{按键与手势模块}
四键经两级同步和约10ms消抖。KEY0循环原图、风格图、左右对比；KEY1/KEY2下一/上一风格；KEY3按下一次切换UI，长按只翻转一次。KEY1+KEY2长按约1秒进入或退出学习，学习时KEY0录入，KEY1/KEY2选择标签。
\sig{gesture_snapshot}从原图左侧$x=0\ldots639$、$y=40\ldots679$旁路采样，每10像素取一处，得到64×64 RGB565、8192B快照。REQUEST后等待新SOF、4096样本及好EOF才ready，CPU多次扫描期间不覆写，读完再ACK。seq、epoch和0.6秒源龄拒绝重复及过期数据。
当前分类需要3帧空背景和每类4份本人模板。补偿整体亮度后保留最大连通前景，规范化并提取64维HOG、64维灰度、64维轮廓，比较镜像距离并拒绝不明确结果。它不设固定R大于G的肤色门，但仍不是通用手部检测器。
约333ms一张快照，三张多数投票形成约1秒决定；至少5次同类新决定且完整4秒才原子提交。RAW确认后自动切STYLE，SPLIT确认后保持对比。STYLE暂停识别，返回RAW/SPLIT可继续。同一手势持续保持不重复触发，不同数字重新满足确认窗口。

\begin{table}[H]\centering\small
\caption{用户手形与十风格顺序，原风格编号依次为0、2、3、4、34、7、9、15、31、40}
\begin{tabularx}{\linewidth}{@{}clX@{}}\toprule 数字 & 风格 & 手形\\\midrule
0 & 梵高旋涡厚涂 & 五指弯曲围成O形\\
1 & 浮世绘海浪 & 食指\\
2 & 彩色剪纸 & 食指、中指\\
3 & 复古漫画 & 中指、无名指、小指\\
4 & 竹林水墨 & 除拇指外四指\\
5 & 阿什维尔抽象 & 五指张开\\
6 & 形体对比 & 拇指、小指\\
7 & 蒙德里安几何 & 拇指、食指、小指\\
8 & 彩色马赛克 & 拇指、食指\\
9 & 雨夜厚涂 & 食指弯钩\\\bottomrule
\end{tabularx}
\end{table}
模板断电丢失。通用免本人录入模型已试训，但人员留出结果和第7/9类不足，未装入当前bit，不能称已经通用识别0--9。


\blocktitle{快照与界面状态的协同}
快照状态顺序为空闲→pending→active→ready→ACK。CPU请求仅预约下一帧；SOF后开始采样，4096点及完整好EOF后才ready。CPU可多次读取这块固定内存，FPGA不得在读取期间覆写；最后一次读取后重新检查seq、epoch、ready和年龄，再ACK释放。

640×640青框不是CNN裁剪窗口。它仅用于手势，实际采样点为$(4+10n,45+10m)$，$n,m=0\ldots63$，分布在原图左半区域。选择横坐标4是因为输入一拍两个像素、计数器只遍历左侧偶坐标；若误写5，采样条件永远不满足。主视频仍然1280×720，手势快照只占8192B，且不额外占用DDR读端口。

视图、手动风格或学习事件使epoch更新，旧快照立即失效。识别资格取决于实际已生效的原图/对比界面，不能用尚在排队的请求模式代替。KEY3只改变文字及青框可见性，隐藏UI后符合条件的手势采样继续运行。

\blocktitle{前景、规范化和分类距离}
背景由三张空框图估计；先补偿整体亮度变化，再根据背景差异和噪声形成二值前景。邻域清理去掉散点，8连通标记保留最大区域；其包围框归一化为64×64，并以5\%/95\%亮度分位数调整范围，减少单个反光点对对比度的影响。这里没有固定肤色门，偏绿图像不会仅因R小于G而被全部丢弃，但移动背景、遮挡和严重弱光仍可能导致前景错误。

特征由三部分组成：4×4空间格、每格4方向的无向HOG共64维；8×8网格灰度均值64维；8×8网格前景密度64维。模板与样本的加权距离为
\[
D=\frac{D_{\rm HOG}+\lfloor D_{\rm gray}/2\rfloor+2D_{\rm mask}}{224},
\]
各项先累加对应维度的绝对差，最后按C整数规则取整。再比较水平镜像后的距离，取较小值。每类取四模板中的最小距离，在十类中选最优及次优；拒绝阈值为该类模板最大两两距离的两倍加8，限制在14--40之间。最优距离超过阈值，或次优差距小于$\max(\lfloor D_{\rm best}/8\rfloor,2)$时，不输出确定数字。

该方法把“不明确”与“无手”分开，数字0则用有效code 1表示，避免0被误当成未检测。个人模板可适应用户定义的O形0、后三指3与拇食小7，代价是需要本人录入且模板断电丢失。

\blocktitle{四秒确认与重复识别}
三张独立新快照中至少两张一致才形成一次候选决定；无多数或歧义不直接增加保持时间。软件同时检查两项条件：连续至少五次同类有效决定，以及首末决定间已过400000000个100MHz UI tick。第一票记为$t_0$，后续约$t_0+1,2,3,4$秒的五票才能证明四个完整秒；不是四票就等于四秒。

请求年龄超过0.6秒、报告间隔超过1.6秒、重复seq、epoch变化或类别中断均清除资格。通过时先写epoch token，再在APB同一接受边沿验证token与风格ID并提交；此时若按键/模式已改变，请求被拒绝。P表示时间进度，并不独立证明提交被接受；诊断需结合原子提交结果与前台风格S判断。

确认后同一手势保持不会连续触发。持续无手达到重武装条件，或检测到不同有效数字后重新满足完整四秒，可再次切换；切STYLE时停止识别，按KEY0返回RAW/SPLIT后恢复。该机制减少抖动与“一次后永久锁住”的状态问题，但不能纠正持续且稳定的错误分类。

\subsubsection{PC训练与诊断工具}
PC离线完成教师引导、QAT、参数导出、整数参考及报告图生成，运行时不为摄像头链路提供推理。CAM7串口为115200、8N1、33个HEX32字段，输出采集、作业、快照与错误状态；它们是观察值，不代替原子握手。


\blocktitle{诊断字段的使用方法}
CAM7一条记录含33个32位十六进制字段。前部保留包、帧、引擎与DDR错误，后部24--31描述快照seq、CPU报告数、分类reason、前景量/阈值、距离/间隔、RGB均值、CPU耗时及epoch/ready等。最后mode字记录请求、加载、作业与前台身份，这四者在切换过程中允许短时不同，不应把requested直接标成“屏幕已切换”。

排查时按数据流顺序定位：C不增长先看传感器/CSI/好帧发布；C增长而I不增长看源池、参数有效、START与错误IRQ；I增长但屏幕不更新看commit及reader欠载；手势报告不增长看ready/seq/enable，报告增长而无候选看前景和距离拒绝。不同寄存器分次更新，串口行不能代替算法内部原子快照；不能据某一行的P和另一个时刻的S拼出一次确认事件。
